\section{From concentration to a limiting location}\label{sec:convergence}

The preceding estimates control fluctuations at one size.  We now rule
out drift of the centers as the number of variables grows.  The argument
has two steps.  An auxiliary Poisson formula compares the satisfiability
probabilities at a combined size and at two smaller sizes.  Concentration
turns this probability comparison into inequalities between centers, whose
errors can then be summed along a balanced tree.

Recall that \(H_n=\inf\{m\ge1:F_{n,m}\notin\sat\}\) is the first
unsatisfiable index in the proper-clause process.  To state the comparison
argument independently of the cap used above, fix \(\lambda>0\) and define
\begin{equation}\label{eq:general-capped-center}
\begin{gathered}
 M_n^{(\lambda)}=\lfloor\lambda n\rfloor,\qquad
 V_n^{(\lambda)}=\min\{H_n-1,M_n^{(\lambda)}\},\\
 x_n^{(\lambda)}=\frac{\EE V_n^{(\lambda)}}{n}.
\end{gathered}
\end{equation}
Thus, for every integer \(0\le m\le M_n^{(\lambda)}\),
\begin{equation}\label{eq:general-cap-survival}
 P_n(m)=\PP\bigl(V_n^{(\lambda)}\ge m\bigr).
\end{equation}
This identity includes \(m=M_n^{(\lambda)}\).  The superscript keeps
these statistics distinct from the particular \(V_n\) and \(\mu_n\)
used in the preceding sections.

\Needspace{24\baselineskip}
\begin{samepage}
\begin{theorem}[Concentration implies convergence]
\label{thm:concentration-to-convergence}
Fix an integer \(k\ge3\) and \(\lambda>0\).  Suppose that the proper
independent uniformly signed \(k\)-clause process has constants
\[
 0<a\le b<\lambda,\qquad 0<\delta<\frac12,\qquad
 \eta>0,\qquad 0\le K<\infty
\]
such that, for all sufficiently large \(n\),
\begin{equation}\label{eq:convergence-input}
\begin{gathered}
 a\le x_n^{(\lambda)}\le b,\\
 \PP\left(\left|V_n^{(\lambda)}-nx_n^{(\lambda)}\right|
             \ge n^{1-\delta}\right)\le K n^{-\eta}.
\end{gathered}
\end{equation}
Then \(x_n^{(\lambda)}\) converges to some \(\alpha\in[a,b]\), and
\[
 \lim_{n\to\infty}P_n(\lfloor cn\rfloor)=
 \begin{cases}
  1,&0\le c<\alpha,\\
  0,&c>\alpha.
 \end{cases}
\]
Moreover, for all sufficiently large \(r,s\),
\begin{align}
 x_{r+s}^{(\lambda)}
 &\ge\min\{x_r^{(\lambda)},x_s^{(\lambda)}\}
       -4\min\{r,s\}^{-\delta},\label{eq:convergence-sum}\\
 x_{s+1}^{(\lambda)}
 &\ge x_s^{(\lambda)}-4s^{-\delta}.
 \label{eq:convergence-neighbor}
\end{align}
The first comparison holds without a restriction on \(r/s\).
\end{theorem}
\end{samepage}

The input is a tail estimate about a finite-size mean.  The conclusion
gives convergence of that mean, while the quantitative comparison is
one-sided.  We begin with the probability inequality that supplies it.

\subsection{An auxiliary product inequality}

For every positive integer \(n\) and \(c\ge0\), let \(s_n(c)\) be the
satisfiability probability of the following auxiliary formula.  Its number
of clauses is \(\Pois(cn)\).  Within each clause, the \(k\) variable indices
are chosen independently and uniformly from \(\{1,\ldots,n\}\), and their
signs are independent and fair.  The clauses are independent.  The
resulting literals are disjoined in the usual Boolean sense, so repeated
variables and tautologies are allowed.  This definition applies at every
positive size, including \(n=1\).

The following calculation is a Poisson form of the frozen-variable
interpolation of Bayati, Gamarnik, and Tetali
\cite[Section~4, Proposition~2]{BGT}.  Independence of the endpoint
Poisson families gives the product that we need.

\begin{lemma}[Auxiliary product inequality]\label{lem:auxiliary-product}
For every pair of positive integers \(r,s\) and every \(c\ge0\),
\begin{equation}\label{eq:auxiliary-product}
 s_{r+s}(c)\ge s_r(c)s_s(c).
\end{equation}
\end{lemma}

\begin{proof}
Let \(N=r+s\), and partition the variables into parts of sizes
\(n_1=r\) and \(n_2=s\).  For \(0\le t\le c\), take three independent
Poisson families: a mean \(Nt\) family of auxiliary clauses on all
\(N\) variables and, for \(i=1,2\), a mean \(n_i(c-t)\) family of
auxiliary clauses confined to part \(i\).  Write \(G_t\) for their
conjunction and \(f(t)=\PP(G_t\in\sat)\).

We use the elementary add-one formula for a marked Poisson family.
If a bounded observable has conditional expectations \(a_j\) given
exactly \(j\) independent marks, then its expectation at mean \(v\) is
\(\sum_{j\ge0}e^{-v}v^ja_j/j!\), whose derivative is
\(\sum_{j\ge0}e^{-v}v^j(a_{j+1}-a_j)/j!\).  Both series converge
uniformly on compact intervals of \(v\), so the formula also applies
after conditioning on the other two Poisson families.

For a satisfiable background \(G_t\), let \(b_i\) be the number of
variables in part \(i\) that are forced by the entire background: every
satisfying assignment gives each such variable the same value.  Set
\(b_i=0\) on unsatisfiable backgrounds.  For a satisfiable background,
a new auxiliary clause destroys satisfiability exactly when each of its literals
is false in every satisfying assignment.  Each literal must therefore
oppose the value of a forced variable.  Conversely, that condition makes
the whole new clause false in every solution.  The independent signed
draws give killing probabilities
\[
 \left(\frac{b_1+b_2}{2N}\right)^k
 \quad\text{and}\quad
 \left(\frac{b_i}{2n_i}\right)^k
\]
for a global clause and a clause confined to part \(i\), respectively.
This remains exact when an index is repeated: each independent signed
draw must still be one of the literals false in every solution.

An unsatisfiable background remains unsatisfiable after any clause is
added.  The add-one formula, applied to the three means, therefore gives
for \(0<t<c\)
\begin{align*}
 f'(t)
 &=N\EE\left[\ind_{\{G_t\in\sat\}}
   \left\{\sum_{i=1}^{2}\frac{n_i}{N}
       \left(\frac{b_i}{2n_i}\right)^k
       -\left(\frac{b_1+b_2}{2N}\right)^k\right\}\right]\\
 &\ge0.
\end{align*}
The last inequality is convexity of \(z\mapsto z^k\) on
\([0,\infty)\), since \((b_1+b_2)/(2N)\) is the weighted average of
\(b_i/(2n_i)\) with weights \(n_i/N\).  At \(t=0\) the two restricted
formulas are independent, so \(f(0)=s_r(c)s_s(c)\); at \(t=c\), only
the global formula remains, so \(f(c)=s_N(c)\).  Continuity at the
endpoints proves the assertion for \(c>0\), and \(c=0\) is immediate.
\end{proof}

\subsection{Transfer to proper clauses}

For the rest of the proof of
Theorem~\ref{thm:concentration-to-convergence}, assume its hypotheses
and keep \(k,\lambda,a,b,\delta,\eta,K\) fixed.  Choose any density
ceiling \(D>b\).  The elementary sandwich below does not use the
concentration hypothesis; its tail consequences also use
\eqref{eq:convergence-input}.  The ceiling is independent of the cap
parameter \(\lambda\): the cap locates the survival events in
\eqref{eq:general-cap-survival}, whereas \(D\) bounds the expected number
of clauses with internal repetitions in the auxiliary formula.  Put
\begin{equation}\label{eq:transfer-parameters}
 R=D\binom{k}{2},\qquad
 \gamma=\min\{\eta,1-2\delta\}>0,\qquad
 \kappa=\frac{e^{-R}}2.
\end{equation}

\begin{samepage}
\begin{lemma}[Proper-clause transfer]\label{lem:proper-transfer}
For \(n\ge k\) set \(\theta_n=(n)_k/n^k\), where
\((n)_k=n(n-1)\cdots(n-k+1)\).  If \(0\le c\le D\) and
\(J\sim\Pois(cn\theta_n)\), then
\begin{equation}\label{eq:proper-sandwich}
\begin{gathered}
 e^{-R}\EE P_n(J)\le s_n(c)\le\EE P_n(J),\\
 cn-R\le\EE J\le cn,\qquad \Var(J)\le Dn.
\end{gathered}
\end{equation}
If \eqref{eq:convergence-input} holds, then there are constants
\(0\le K'<\infty\) and an integer \(n_1\ge k\) such that for every \(n\ge n_1\) and
every \(c\in[0,D]\),
\begin{align}
 c\le x_n^{(\lambda)}-2n^{-\delta}
 &\quad\Longrightarrow\quad s_n(c)\ge\kappa,
 \label{eq:proper-transfer-low}\\
 c\ge x_n^{(\lambda)}+2n^{-\delta}
 &\quad\Longrightarrow\quad s_n(c)\le K'n^{-\gamma}.
 \label{eq:proper-transfer-high}
\end{align}
The constants \(K'\) and \(n_1\) may depend on the fixed parameters,
but not on \(c\) or \(n\).
\end{lemma}
\end{samepage}

\begin{proof}
An auxiliary clause has distinct indices with probability \(\theta_n\).
Poisson thinning separates the clauses into a family of \(J\) clauses
with distinct indices and a family of \(Q\) clauses with an internal
repetition, where
\[
 J\sim\Pois(cn\theta_n),\qquad Q\sim\Pois(cn(1-\theta_n)).
\]
These counts are independent.  Indeed, their joint probability-generating
function is
\[
 \EE[z^Jw^Q]
 =\exp\bigl(cn[\theta_n z+(1-\theta_n)w-1]\bigr)
 =e^{cn\theta_n(z-1)}e^{cn(1-\theta_n)(w-1)}.
\]
The same independent-mark construction shows that, conditional on these
counts, the marks in the two families are independent with their
respective conditional laws.  In the first family, every unordered set
of \(k\) distinct variables has \(k!\) equally likely orders and the
signs remain independent and fair.  Conditional on \(J\), this family
therefore has exactly the law \(F_{n,J}\).

A union bound over the \(\binom{k}{2}\) pairs of variable positions gives
\[
 cn(1-\theta_n)\le c\binom{k}{2}\le R.
\]
Consequently \(cn-R\le\EE J\le cn\) and \(\Var(J)=\EE J\le Dn\).
Deleting the repeated-variable clauses can only help satisfiability.
Requiring \(Q=0\), which is independent of the proper family and has
probability \(e^{-cn(1-\theta_n)}\ge e^{-R}\), gives the lower bound.
This proves \eqref{eq:proper-sandwich}.

Write \(h=n^{1-\delta}\).  The concentration input and the survival
identity give, for all sufficiently large \(n\) and integer \(j\),
\begin{align*}
 P_n(j)&\ge1-Kn^{-\eta}
 &&\text{if }0\le j\le nx_n^{(\lambda)}-h,\\
 P_n(j)&\le Kn^{-\eta}
 &&\text{if }j\ge nx_n^{(\lambda)}+h.
\end{align*}
For the first estimate the indicated counts are below the cap because
\(x_n^{(\lambda)}\le b<\lambda\).  For the second, first use
\eqref{eq:general-cap-survival} up to the cap and then use monotonicity
beyond it at \(j_0=\lceil nx_n^{(\lambda)}+h\rceil\).
The inequality \(j_0\le M_n^{(\lambda)}\) holds for all sufficiently
large \(n\): the gap \((\lambda-b)n\) dominates \(h+2\).  This also
accounts for both integer roundings.

If \(c\le x_n^{(\lambda)}-2n^{-\delta}\), then
\(\EE J\le nx_n^{(\lambda)}-2h\).  Chebyshev's inequality yields
\[
 \PP\bigl(J>nx_n^{(\lambda)}-h\bigr)
 \le\frac{Dn}{h^2}=D n^{-(1-2\delta)}.
\]
\begin{samepage}
The lower bound in \eqref{eq:proper-sandwich} is thus at least
\[
 e^{-R}\bigl(1-Kn^{-\eta}-D n^{-(1-2\delta)}\bigr),
\]
which is at least \(\kappa\) for every sufficiently large \(n\).
\end{samepage}

If \(c\ge x_n^{(\lambda)}+2n^{-\delta}\), then
\(\EE J\ge nx_n^{(\lambda)}+2h-R\).  For large \(n\), \(h>R\), and
\[
 \PP\bigl(J<nx_n^{(\lambda)}+h\bigr)
 \le\frac{Dn}{(h-R)^2}=O\bigl(n^{-(1-2\delta)}\bigr).
\]
The upper bound in \eqref{eq:proper-sandwich} is at most this probability
plus \(Kn^{-\eta}\), proving \eqref{eq:proper-transfer-high} with
\(\gamma=\min\{\eta,1-2\delta\}\).  All estimates used the common
bounds \(R\) and \(Dn\), so they are uniform on \([0,D]\).
\end{proof}

\subsection{Comparing centers at different sizes}

The lower transfer bound \eqref{eq:proper-transfer-low} stays positive,
whereas the upper bound \eqref{eq:proper-transfer-high} tends to zero
uniformly.  We use this separation with the product inequality
\eqref{eq:auxiliary-product} to compare the centers.

Let \(u=\min\{r,s\}\), \(N=r+s\), and set
\[
 c=\min\{x_r^{(\lambda)},x_s^{(\lambda)}\}-2u^{-\delta}.
\]
For all sufficiently large \(u\), the center bounds in
\eqref{eq:convergence-input} put \(c\) in \([0,D]\).  Since
\(u^{-\delta}\ge r^{-\delta},s^{-\delta}\), the lower estimate in
Lemma~\ref{lem:proper-transfer} applies at both sizes.  Hence
Lemma~\ref{lem:auxiliary-product} gives \(s_N(c)\ge\kappa^2\).
For all sufficiently large \(N\), the upper estimate in
Lemma~\ref{lem:proper-transfer} is smaller than \(\kappa^2\); it rules
out \(c\ge x_N^{(\lambda)}+2N^{-\delta}\).  Therefore
\[
 x_N^{(\lambda)}>c-2N^{-\delta}
 \ge\min\{x_r^{(\lambda)},x_s^{(\lambda)}\}-4u^{-\delta}.
\]
This proves \eqref{eq:convergence-sum} for every pair of sufficiently
large sizes.  No ratio between them was used.

To reach every sufficiently large size from a fixed block size, we will
use blocks of sizes \(s\) and \(s+1\).  We therefore also need a
neighboring-size comparison.  Take
\(c=x_s^{(\lambda)}-2s^{-\delta}\), which again lies in \([0,D]\)
for large \(s\).  In the size-one auxiliary model, the empty-formula
event gives \(s_1(c)\ge e^{-c}\ge e^{-D}\).  Consequently
\[
 s_{s+1}(c)\ge s_s(c)s_1(c)\ge\kappa e^{-D}.
\]
For large \(s\), the upper transfer estimate at size \(s+1\) is smaller
than this fixed positive constant.  It follows that
\[
 x_{s+1}^{(\lambda)}>c-2(s+1)^{-\delta}
 \ge x_s^{(\lambda)}-4s^{-\delta},
\]
which proves \eqref{eq:convergence-neighbor}.

The tree argument below uses the sum comparison only when its child
sizes are comparable.

\subsection{Summing the errors along a balanced tree}

The principle that errors summable over geometrically increasing sizes
cannot sustain macroscopic drift parallels the approximate-superadditivity
argument of Abbe and Montanari \cite[arXiv version, Lemma~8]{AM}.
Here the operation is a minimum of normalized quantities; the following
lemma proves the required form directly.

\begin{samepage}
\begin{lemma}[Balanced comparisons]\label{lem:balanced-comparison}
Let \(n_0\ge1\) be an integer and let \((u_n)_{n\ge n_0}\) be a bounded
real sequence.  Suppose that \(0\le A<\infty\), \(\beta>0\), and that,
for all \(r,s\ge n_0\) with \(1/3\le r/s\le3\),
\[
 u_{r+s}\ge\min\{u_r,u_s\}-A\min\{r,s\}^{-\beta}.
\]
Suppose also that, for every \(s\ge n_0\),
\[
 u_{s+1}\ge u_s-As^{-\beta}.
\]
Then \(u_n\) converges.  More precisely, put
\[
 B_\beta=\frac{(3/2)^\beta}{1-(2/3)^\beta}.
\]
For every integer \(s\ge\max\{n_0,2\}\) and every integer \(n\ge s^2\),
\begin{equation}\label{eq:balanced-comparison-bound}
 u_n\ge u_s-A(1+B_\beta)s^{-\beta}.
\end{equation}
\end{lemma}
\end{samepage}

\begin{proof}
Fix an integer \(s\ge\max\{n_0,2\}\).  For \(n\ge s^2\),
write \(n=qs+t\), where \(q=\lfloor n/s\rfloor\ge s\) and \(0\le t<s\).
Thus \(n\) is a sum of \(q-t\) blocks of size \(s\) and \(t\) blocks
of size \(s+1\).  Both block counts are nonnegative, and the neighbor
inequality puts the value of every block at least
\(u_s-As^{-\beta}\).

Arrange these \(q\) blocks in a binary tree by dividing each group of
\(a\ge2\) blocks into groups of \(\lfloor a/2\rfloor\) and
\(\lceil a/2\rceil\) blocks.  Each child contains at least \(a/3\)
blocks, so its total size is at least \(sa/3\).  Moreover, the ratio of
the larger child's total size to the smaller child's is at most
\[
 \frac{\lceil a/2\rceil(s+1)}{\lfloor a/2\rfloor s}
 \le2(1+1/s)\le3.
\]
The sum inequality is therefore valid at every internal node, with
error at most \(A(sa/3)^{-\beta}\).

Iterating the minimum inequality from the leaves to the root gives a
lower bound by the least leaf value minus the largest sum of errors on
a leaf-to-root path.  If \(a_1,a_2,\ldots\) are the block counts at the
successive internal nodes of such a path, then \(a_1\ge2\) and
\(a_{j+1}\ge(3/2)a_j\), because a child contains at most two thirds of
its parent's blocks.  Every path error is consequently at most
\[
 As^{-\beta}(3/2)^\beta
   \sum_{j=0}^{\infty}(2/3)^{j\beta}
 =AB_\beta s^{-\beta}.
\]
This proves \eqref{eq:balanced-comparison-bound}.  First let
\(n\to\infty\) with \(s\) fixed to obtain
\[
 \liminf_{n\to\infty}u_n
 \ge u_s-A(1+B_\beta)s^{-\beta}.
\]
Taking the limsup as \(s\to\infty\) gives
\(\liminf_n u_n\ge\limsup_n u_n\).  Boundedness makes both quantities
finite, and hence \(u_n\) converges.
\end{proof}

\begin{proof}[Completion of Theorem~\ref{thm:concentration-to-convergence}]
Apply Lemma~\ref{lem:balanced-comparison} to
\eqref{eq:convergence-sum}--\eqref{eq:convergence-neighbor}, with
\(A=4\) and \(\beta=\delta\).  It gives convergence, and the center
bounds place the limit \(\alpha\) in \([a,b]\).

Fix \(0\le c<\alpha\).  Since \(x_n^{(\lambda)}\to\alpha\) and
\(n^{-\delta}\to0\), for all sufficiently large \(n\),
\(\lfloor cn\rfloor\le nx_n^{(\lambda)}-n^{1-\delta}\).  This count
lies below the cap, and \eqref{eq:convergence-input} together with
\eqref{eq:general-cap-survival} gives
\(P_n(\lfloor cn\rfloor)\ge1-Kn^{-\eta}\to1\).

Now fix \(c>\alpha\), and choose a real \(d\) with
\(\alpha<d<\min\{c,\lambda\}\), which is possible because
\(\alpha\le b<\lambda\).  For large \(n\), the integer
\(\lfloor dn\rfloor\) lies below the cap and is at least
\(nx_n^{(\lambda)}+n^{1-\delta}\).  The same two estimates give
\(P_n(\lfloor dn\rfloor)\le Kn^{-\eta}\to0\).
Monotonicity then yields
\(P_n(\lfloor cn\rfloor)\le P_n(\lfloor dn\rfloor)\to0\).
\end{proof}

\subsection{Application to every fixed clause size}

We finish by checking the inputs for the statistic used in this paper.
This also identifies the repetition constant furnished by the general
transfer lemma.

\begin{proof}[Proof of Theorem~\ref{thm:main}]
Fix \(k\ge3\).  Recall \(L_k=2^{k+1}\),
\(V_n=\min\{H_n-1,L_kn\}\), and \(\mu_n=\EE V_n/n\).
Since \(L_k\) is an integer, the definitions in
\eqref{eq:general-capped-center} with \(\lambda=L_k\) give exactly
\(V_n^{(L_k)}=V_n\) and \(x_n^{(L_k)}=\mu_n\).

Proposition~\ref{prop:concentration} supplies the tail estimate in
\eqref{eq:convergence-input} with
\[
 \delta=\delta_k=\frac{k-2}{4k},\qquad
 \eta=\eta_k=\frac{k-2}{2k},
\]
and some \(K=K_k<\infty\).  These exponents satisfy
\(0<\delta<1/2\) and \(\eta>0\).  Lemma~\ref{lem:center-bounds}
supplies the center bounds with
\(a=a_0\) and \(b=2^k+1<L_k\).

For the transfer in this application choose \(D=L_k>b\).  Then
\[
 R=L_k\binom{k}{2},\qquad
 1-2\delta=\frac{k+2}{2k}>\eta,
 \qquad \gamma=\eta,\qquad \kappa=\frac{e^{-R}}2.
\]
Thus all hypotheses of Theorem~\ref{thm:concentration-to-convergence}
hold with the stated cap, center bounds, and exponents.  It follows that
\(\mu_n\to\alpha_k\in[a_0,2^k+1]\) and that the proper-clause
satisfiability probabilities tend to one for \(0\le c<\alpha_k\)
and to zero for \(c>\alpha_k\), as claimed.
\end{proof}
